\documentclass[12pt,a4paper]{article}
\usepackage[T1]{fontenc}
\usepackage[ngerman]{babel}
\usepackage{lmodern}
\usepackage{csquotes}
\usepackage{amsmath,amssymb}
\usepackage{geometry}
\geometry{a4paper, margin=2.5cm}
\usepackage{parskip}
\usepackage{xcolor}
\usepackage[colorlinks=true, linkcolor=blue!60!black, urlcolor=blue!60!black, citecolor=blue!60!black]{hyperref}
\usepackage[backend=biber, style=numeric]{biblatex}
\usepackage[toc, acronym, nonumberlist, style=long]{glossaries}
\makeglossaries
\sloppy
\begin{filecontents*}{\jobname.bib}
@misc{bsi2023smtpsmuggling,
  author = {{Bundesamt für Sicherheit in der Informationstechnik}},
  title = {SMTP Smuggling ermöglicht Social Engineering per E-Mail (CSW-Nr. 2023-292569-1032)},
  year = {2023}
}
\end{filecontents*}

\addbibresource{\jobname.bib}

\newacronym{mfa}{MFA}{Multi-Faktor-Authentisierung}

\newglossaryentry{phishing}{
  name=Phishing,
  description={Eine Form des Social Engineering, bei der Angreifer versuchen, über gefälschte E-Mails, Nachrichten oder Webseiten an vertrauliche Daten zu gelangen.}
}

\newglossaryentry{socialengineering}{
  name=Social Engineering,
  description={Die Manipulation von Personen durch psychologische Tricks, um an vertrauliche Informationen zu gelangen.}
}

\newglossaryentry{spearphishing}{
  name=Spear-Phishing,
  description={Ein gezielter Phishing-Angriff, der auf eine spezifische Person oder bestimmte Mitarbeiter zugeschnitten ist.}
}

\newglossaryentry{quishing}{
  name=Quishing,
  description={Eine Phishing-Methode, bei der manipulierte QR-Codes eingesetzt werden, um Opfer auf gefälschte Webseiten zu leiten.}
}

\begin{document}

% --- TITELBLATT ---
\begin{titlepage}
    \centering
    \vspace*{2cm}
    
    {\Huge \textbf{Skript: Phishing} \par}
    \vspace{1.5cm}
    {\Large \textbf{Kantonsschule im Lee} \par}
    \vspace{2cm}
    
    \rule{\linewidth}{0.5mm}
    \vspace{2cm}
    
    {\large \textbf{Thema:} Phishing, Social Engineering \& Best Practices \par}
    \vspace{0.5cm}
    {\large \textbf{Fach:} Informatik \par}
    \vspace{0.5cm}
    {\large \textbf{Autor:} Laurin Ott \par}
    
    \vfill
    
    {\large \textbf{Datum:} \today \par}
\end{titlepage}

% --- INHALTSVERZEICHNIS ---
\newpage
\tableofcontents
\newpage

\section{Einleitung}
In der digitalen Welt bilden nicht nur technische Barrieren wie Passwörter die Abwehrlinie, sondern vor allem der Mensch selbst. Eines der erfolgreichsten und am häufigsten genutzten Einfallstore für Cyberkriminelle ist \gls{phishing}. Dabei wird gezielt die menschliche Psyche ausgenutzt, um an vertrauliche Daten zu gelangen oder Schadsoftware zu verbreiten \cite{bsi2023smtpsmuggling}.

\section{Grundlegende Begriffe und Angriffsarten}
\gls{phishing} gehört zu den klassischen \gls{socialengineering}-Angriffen. Statt komplexe technische Systeme direkt zu hacken, manipulieren Angreifer gezielt Personen:
\begin{itemize}
    \item \textbf{Klassisches Phishing:} Täuschend echte E-Mails, SMS oder Nachrichten, die vorgeben, von vertrauenswürdigen Institutionen (wie Banken, Streaming-Diensten oder Behörden) zu stammen.
    \item \textbf{Spear-Phishing:} Ein gezielter Angriff (\gls{spearphishing}), der auf eine spezifische Person oder bestimmte Mitarbeiter zugeschnitten ist, um Vertrauen zu erwecken.
    \item \textbf{Vishing \& Smishing:} Betrugsversuche über Telefonanrufe (\textit{Voice}) oder Kurznachrichten (\textit{SMS}).
    \item \textbf{Quishing:} Das Platzieren von manipulierten QR-Codes (z. B. auf Plakaten oder Parkuhren), die zu gefälschten Webseiten führen.
\end{itemize}

\section{Psychologische Faktoren und Wirkungsweise}
Im Gegensatz zu mathematischen Angriffen zielt \gls{phishing} auf menschliche Schwachstellen wie Neugier, Hilfsbereitschaft, Gier oder vor allem \textbf{Angst und Zeitdruck} ab. Typische Maschen beinhalten drohende Konsequenzen wie angebliche Kontosperrungen, dringende Rechnungen oder verpasste Gewinne, damit das Opfer unüberlegt handelt.

\section{Aktuelle Richtlinien und Best Practices}
Um sich vor \gls{phishing}-Angriffen zu schützen, gelten im Rahmen der IT-Sicherheit folgende zentrale Empfehlungen:

\begin{enumerate}
    \item \textbf{Kritische Hürden bewahren:} Niemals unüberlegt auf Links klicken oder Anhänge von unbekannten Absendern öffnen.
    \item \textbf{Absender und URLs genau prüfen:} Auf Unstimmigkeiten in E-Mail-Adressen (z. B. Buchstabendreher in Domains) achten.
    \item \textbf{Multi-Faktor-Authentisierung (MFA) nutzen:} Selbst wenn Anmeldedaten abgegriffen werden, verhindert ein zusätzlicher Faktor (wie eine Bestätigung auf dem Smartphone) unbefugten Zugriff.
    \item \textbf{Keine sensiblen Daten preisgeben:} Seriöse Institutionen fragen niemals per E-Mail, SMS oder Telefon nach Passwörtern oder PINs.
\end{enumerate}

\printbibliography

\printglossary[type=\acronymtype, title=Akronyme]
\printglossary[title=Glossar]

\end{document}